\chapter{Discussion}
\label{ch:discussion}

Returning to the problem statement (Chapter~\ref{ch:intro}): this project set out to
determine whether an array of small apertures can substitute for a corrective lens over
a useful field of view, and to characterize the trade-offs involved. The evidence
assembled here supports a qualified yes, with the trade-offs now quantified rather than
asserted.

\textbf{Single-aperture optics} (Chapter~\ref{ch:single}) establish that the stenopeic
effect is real and substantial (Section~\ref{sec:exp3}: a $15$--$18\times$ reduction in
defocus-growth rate, measured photographically, not merely predicted), and that an
exact wave-optics treatment is necessary, not a convenience: the closed-form RSS
approximation this project also implements diverges from the exact result by 64.8\% on
average, and a specific naive-FFT implementation trap (Section~\ref{sec:bug-fft}) would
have silently inverted the physics of myopic defocus entirely had it shipped
uncorrected.

\textbf{Array-level optics} (Chapter~\ref{ch:array-theory}) show that extending a single
aperture to an array introduces two new failure modes -- ghosting and blinding -- with
no single-aperture analogue, and that a naive per-aperture margin criterion for
ghosting is not merely imprecise but was shown, via the Barten CSF correction
(Section~\ref{sec:barten}), to reclassify the \emph{entire} tested dense-pitch range
from ``safe'' to ``danger zone.'' The harmonic-series extension
(Section~\ref{sec:harmonic}) further shows that even a corrected fundamental-frequency
analysis can underestimate real risk by over $4\times$ by missing the worst-case
harmonic.

\textbf{The two evaluation models} (Chapters~\ref{ch:model-a}--\ref{ch:model-b}) are not
in conflict: Model~B generalizes Model~A by removing its single-axis assumption, and
the fact that removing that assumption changed the outcome -- surfacing a real
blackout failure mode Model~A's own structural-safety guarantee had missed
(Section~\ref{sec:bug-single-axis}) -- is direct evidence the generalization mattered,
not merely a theoretical nicety. That hexagonal packing wins under the more rigorous
model is explained, not merely observed, by a first-principles symmetry argument
(Section~\ref{sec:hex-intuition}) grounded in the crystallographic restriction theorem,
and demonstrated with real, matched-spacing geometry and real, already-computed safety
data (Figures~\ref{fig:geom-structure}--\ref{fig:geom-consequence}).

\textbf{External validation} (Chapter~\ref{ch:model-comparison}) is this project's most
honestly limited result: the acuity model, in its raw form, fails a formal $R^2>0$
significance bar against every real clinical comparison available, both before and
after a genuine, verified optics fix that measurably narrowed (but did not close) the
gap. This paper does not claim the model is externally validated; it claims the model
is \emph{physically grounded}, internally self-consistent (Section~\ref{sec:asm-validation}),
improved measurably when a real defect was corrected (Section~\ref{sec:bug-shadowing}),
and its remaining discrepancy against real data is reported with the exact statistics,
not glossed over.
